% ═══════════════════════════════════════════════════════════════
% MT-01a: Minitest Begrüßungen
% Spanisch Klasse 7 - Sequenz 1: Empezamos
% ═══════════════════════════════════════════════════════════════
% Dauer: 10 Minuten | Punkte: 26 BE
% ═══════════════════════════════════════════════════════════════

\documentclass[11pt, a4paper]{article}

\usepackage[utf8]{inputenc}
\usepackage[T1]{fontenc}
\usepackage[ngerman]{babel}
\usepackage{helvet}
\renewcommand{\familydefault}{\sfdefault}

\usepackage[margin=2cm]{geometry}
\usepackage{enumitem}
\usepackage{tabularx}
\usepackage{booktabs}
\usepackage{xcolor}
\usepackage{fancyhdr}
\usepackage{lastpage}
\usepackage{amssymb}

% FARBEN
\definecolor{spanisch}{HTML}{c41e3a}
\definecolor{grau}{HTML}{888888}
\definecolor{hellgrau}{HTML}{f5f5f5}

\pagestyle{fancy}
\fancyhf{}
\renewcommand{\headrulewidth}{0.4pt}
\renewcommand{\footrulewidth}{0.4pt}

\lhead{Spanisch -- Klasse 7}
\chead{\textbf{MT-01a}}
\rhead{\today}

\lfoot{10 Minuten}
\cfoot{}
\rfoot{Seite \thepage\ von \pageref{LastPage}}

\newcommand{\punktebox}[1]{\hfill\fbox{\small \_\_\_/#1}}
\newcommand{\luecke}[1]{\rule{#1}{0.4pt}}
\newcommand{\es}[1]{\textcolor{spanisch}{\textbf{#1}}}

\begin{document}

% ═══════════════════════════════════════════════════════════════
% KOPFBEREICH
% ═══════════════════════════════════════════════════════════════

\begin{center}
    {\Large\bfseries\textcolor{spanisch}{Minitest: Begrüßungen}}
\end{center}

\vspace{5pt}

\noindent
\begin{tabularx}{\textwidth}{|X|X|X|}
\hline
\textbf{Name:} \luecke{4cm} &
\textbf{Klasse:} 7 &
\textbf{Datum:} \luecke{2cm} \\
\hline
\end{tabularx}

\vspace{10pt}

\noindent
\fcolorbox{spanisch}{hellgrau}{%
\parbox{\dimexpr\textwidth-2\fboxsep-2\fboxrule}{%
\textbf{Zeit:} 10 Minuten | \textbf{Punkte:} 26 BE | \textbf{Hilfsmittel:} keine
}}

\vspace{15pt}

% ═══════════════════════════════════════════════════════════════
% AUFGABE 1: Zuordnung (6 BE)
% ═══════════════════════════════════════════════════════════════

\noindent\textbf{\textcolor{spanisch}{Aufgabe 1 -- Begrüßungen zuordnen}} \punktebox{6}

\noindent Verbinde die spanische Begrüßung mit der deutschen Bedeutung. Schreibe den Buchstaben in die Lücke.

\vspace{10pt}

\begin{minipage}[t]{0.48\textwidth}
\begin{enumerate}[label=\arabic*.]
    \item \es{¡Hola!} \luecke{1cm}
    \item \es{¡Buenos días!} \luecke{1cm}
    \item \es{¡Buenas noches!} \luecke{1cm}
    \item \es{¡Adiós!} \luecke{1cm}
    \item \es{¡Hasta luego!} \luecke{1cm}
    \item \es{¡Hasta mañana!} \luecke{1cm}
\end{enumerate}
\end{minipage}
\hfill
\begin{minipage}[t]{0.48\textwidth}
\begin{itemize}[label=]
    \item a) Guten Morgen!
    \item b) Hallo!
    \item c) Guten Abend!
    \item d) Auf Wiedersehen!
    \item e) Bis später!
    \item f) Bis morgen!
\end{itemize}
\end{minipage}

\vspace{15pt}

% ═══════════════════════════════════════════════════════════════
% AUFGABE 2: Tageszeit (6 BE)
% ═══════════════════════════════════════════════════════════════

\noindent\textbf{\textcolor{spanisch}{Aufgabe 2 -- Welche Begrüßung passt?}} \punktebox{6}

\noindent Schreibe die passende Begrüßung: \textit{Buenos días}, \textit{Buenas tardes} oder \textit{Buenas noches}.

\vspace{10pt}

\begin{tabularx}{\textwidth}{|l|X|}
\hline
\textbf{Uhrzeit} & \textbf{Begrüßung} \\
\hline
8:00 Uhr morgens & \luecke{4cm} \\[0.6em]
\hline
15:00 Uhr nachmittags & \luecke{4cm} \\[0.6em]
\hline
21:00 Uhr abends & \luecke{4cm} \\[0.6em]
\hline
10:00 Uhr vormittags & \luecke{4cm} \\[0.6em]
\hline
17:00 Uhr nachmittags & \luecke{4cm} \\[0.6em]
\hline
23:00 Uhr nachts & \luecke{4cm} \\[0.6em]
\hline
\end{tabularx}

\vspace{15pt}

% ═══════════════════════════════════════════════════════════════
% AUFGABE 3: Grammatik (6 BE)
% ═══════════════════════════════════════════════════════════════

\noindent\textbf{\textcolor{spanisch}{Aufgabe 3 -- Buenos oder Buenas?}} \punktebox{6}

\noindent Ergänze die richtige Form: \textit{Buenos} oder \textit{Buenas}.

\vspace{10pt}

\begin{enumerate}[label=\alph*)]
    \item \luecke{2cm} días
    \item \luecke{2cm} tardes
    \item \luecke{2cm} noches
\end{enumerate}

\vspace{10pt}

\noindent \textbf{Begründung (3 BE):} Warum heißt es ``Buenos días'' aber ``Buenas tardes''? Kreuze an.

\vspace{5pt}

\noindent
$\square$ Weil ``día'' maskulin und ``tarde/noche'' feminin ist \\
$\square$ Weil ``día'' kürzer ist \\
$\square$ Weil es morgens ist

\vspace{15pt}

% ═══════════════════════════════════════════════════════════════
% AUFGABE 4: Verabschiedung (4 BE)
% ═══════════════════════════════════════════════════════════════

\noindent\textbf{\textcolor{spanisch}{Aufgabe 4 -- Verabschiedungen ergänzen}} \punktebox{4}

\noindent Ergänze die fehlenden spanischen Wörter.

\vspace{10pt}

\begin{enumerate}[label=\alph*)]
    \item Bis bald! = ¡Hasta \luecke{3cm}!
    \item Bis später! = ¡Hasta \luecke{3cm}!
\end{enumerate}

\vspace{15pt}

% ═══════════════════════════════════════════════════════════════
% AUFGABE 5: Satzzeichen (4 BE)
% ═══════════════════════════════════════════════════════════════

\noindent\textbf{\textcolor{spanisch}{Aufgabe 5 -- Spanische Satzzeichen}} \punktebox{4}

\noindent Was ist die Besonderheit bei spanischen Ausrufe- und Fragezeichen? Kreuze an.

\vspace{5pt}

\noindent
$\square$ Sie stehen nur am Ende des Satzes \\
$\square$ Sie stehen am Anfang UND am Ende des Satzes \\
$\square$ Sie werden nicht verwendet

\vspace{10pt}

\noindent Schreibe ``Hallo, wie geht's?'' auf Spanisch mit korrekten Satzzeichen:

\vspace{5pt}

\noindent \luecke{10cm}

% ═══════════════════════════════════════════════════════════════
% PUNKTEÜBERSICHT
% ═══════════════════════════════════════════════════════════════

\vspace{20pt}
\hrule
\vspace{10pt}

\noindent
\begin{tabularx}{\textwidth}{|X|c|c|c|c|c||c|}
\hline
\textbf{Aufgabe} & 1 & 2 & 3 & 4 & 5 & \textbf{Gesamt} \\
\hline
\textbf{max. Punkte} & 6 & 6 & 6 & 4 & 4 & \textbf{26} \\
\hline
\textbf{erreicht} & & & & & & \\[0.6em]
\hline
\end{tabularx}

\vspace{10pt}

\begin{flushright}
\textbf{Note:} \luecke{2cm}
\end{flushright}

\end{document}
